% =====================================================================
%  GMN2 §3.2.4–3.2.7
%    映射到 Hitchin 系统 / 带基本物质的线性共形箭图 /
%    弯膜与不规则奇点 / 渐近自由情形
%  原文：wkb-paper.tex 2794–3587 行；编号公式 (3.73)–(3.132) 共 60 条；插图 2 幅
%  说明：BRIEF 表把本块行号记作 2794–3263，但那一段只到 §3.2.5 末尾；
%        §3.2.6 起于 3264 行、§3.2.7 起于 3412 行、§3.2.8 起于 3588 行。
%        为覆盖块名所列的四个小节，本文件译至 3587 行止。
%  编号自检：(3.73) 紧接 §3.2.3 的 (3.72)；(3.74)–(3.76) 与 §1 所说
%        「正则奇性的情形见 (3.74)、(3.75)、(3.76)」一致；
%        (3.115)、(3.116) 与 §1 所说「不规则奇性的情形见 (3.115)、(3.116)」一致。
% =====================================================================

\subsubsection{映射到 Hitchin 系统}\label{subsubsec:MapToHitchin}

现在我们可以回到第 3.1 节% TODO 合稿改 \ref{sec:scft}
所描述的 Hitchin 系统了。我们把方程 $F(t,v)=0$ 的解的 $K$ 个分支 $\tilde v(t)$ 解释为 Higgs 场 $\varphi_s$ 的 $K$ 个本征值，因为它们正是 M5-膜的位置。我们固定质心以消去多项式 $F(t,v)$ 中的 $\tilde v^{K-1}$ 项，这一选择保证了 $\varphi_s$ 取值于 $SU(K)$ 而不取值于 $U(K)$。柱面上 Hitchin 系统的谱曲线按定义就是 $\det (\tilde v + \varphi_s) =0$。我们要把这条谱曲线等同于红外 M5-膜所缠绕的曲线 $F(t,v)=0$。唯一要紧的一点是记住：$\varphi_s$ 按一形式变换，即 $\varphi_t dt= \varphi_s ds$。把谱曲线与曲线 $F(t,v)=0$ 等同起来，并作定义 $x := \frac{\tilde v}{t}$，则 $T^* \IC^\times$ 中的谱曲线具有如下形式
\begin{equation}\label{eq:MatchHitchin}
% 原文此处另有一行被注释掉的旧写法（\det(x-\varphi_t)=x^K+\sum \frac{p_i(t)}{t^i\prod (t-t_\alpha)}x^{K-i}=0），未译。
\det (x-\varphi_t) = x^K + \sum_{i=2}^{K} \frac{\tilde p_i(t)}{ \left( t \prod_{\alpha=0}^{n} (t-t_\alpha) \right)^i } x^{K-i} = 0.
\end{equation}
（$\tilde p_i$ 与式 (3.45)% TODO 合稿改 \eqref{eq:swcurve1}
中出现的 $p_i$ 不同，因为我们对 $v$ 作了到 $\tilde v$ 的平移。）如果我们不把 $U(1)$ 自由度除掉，则要等同 $x=v/t$。无论哪种情形，这个方程都应与式 (3.38)% TODO 合稿改 \eqref{eq:speccurve}
等同。注意 SW 微分是 $x\,dt$，正如第 3.1 节所断言的那样。

为了更对称地处理这些奇点，把 $\IC^\times \to \IC \IP^1$ 通过 $t \to [t:1]$ 嵌入进来是有用的。于是我们看出，柱面两端（即 $t=0$ 与 $t=\infty$ 处）的边界条件说的是：$\varphi_t dt$ 有一个单极点，其留数由 $v$ 的渐近值给出。由此我们推出 Higgs 场上的边界条件：
\begin{equation}\label{eq:SimpleAtZero}
\varphi_t dt  \to \frac{dt}{t} {\rm Diag}\{v^{(0)}_1, \dots, v^{(0)}_K
\}\qquad t\to 0
\end{equation}
\begin{equation}\label{eq:SimpleAtInfty}
\varphi_t dt  \to \frac{dt}{t} {\rm Diag}\{v^{(n+1)}_1, \dots, v^{(n+1)}_K
\}\qquad t \to \infty
\end{equation}

此外，在 $t=t_\alpha$ 处还有单极点，其留数形如
\begin{equation}\label{eq:SimplePoleResidue}
{\rm Diag}\{-(K-1) m_\alpha, m_\alpha,m_\alpha,\cdots m_\alpha,m_\alpha \}  ,\qquad \alpha =
0, \dots, n
\end{equation}
这是把 $U(1)$ 自由度除掉之后的情形；若不除掉，则留数为 ${\rm Diag}\{\rho_\alpha, 0,\dots, 0 \}$。回忆 $m_\alpha$ 由式 (3.62)% TODO 合稿改 \eqref{eq:def-m-alpha}
给出。注意 $t=0, \infty$ 处的留数是 $su(K)$ 的一般半单元素，因而对应正则奇点；但当我们考虑 $t=t_\alpha$（$\alpha = 0, \dots, n$）处的留数时，对 $K>2$ 这完全不对，因为此时（对标准单根选取而言）这些留数被除第一个单根以外的所有单根零化。

我们要强调，条件 \eqref{eq:SimpleAtZero}、\eqref{eq:SimpleAtInfty} 与 \eqref{eq:SimplePoleResidue} 的意思是：存在一个局部规范，使得场可以被写成这种形式。

本文主体将集中于 $K=2$ 的情形。此时，按上述方式除掉质心 $U(1)$ 之后，两类奇点之间不再有任何区别。于是 Hitchin 系统成为 $\IC \IP^1$ 上带 $n+3$ 个正则奇点的一般 $SU(2)$ Hitchin 系统。方程 \eqref{eq:MatchHitchin} 等价于
\begin{equation}\label{eq:su2exple}
{\rm tr} \varphi_t^2 = \half \frac{ (p_1(t))^2 - 4 p_2(t)
\prod_{\alpha=0}^n (t-t_\alpha) }{(t \prod_0^n (t-t_\alpha))^2}.
\end{equation}
注意 $\varphi_t dt$ 在 $t = 0$ 与 $t = \infty$ 处、以及在 $n+1$ 个点 $t_a$ 处都有单极点。我们将在第 9 节% TODO 合稿改 \ref{section:Superconformal}
回到这个公式。

\subsubsection{带基本物质的线性共形箭图}
\label{subsubsec:LinConfQuivMatt}

把这一分析推广到带基本物质的更一般线性共形箭图是直截了当的。这是一种含 $SU(k_\alpha)$ 规范群的箭图，在第 $\alpha$ 个节点处有 $d_\alpha = 2 k_\alpha - k_{\alpha-1} - k_{\alpha + 1}$ 个基本场，$\alpha = 1, \dots, n$。这些额外的基本场在膜构造中由 $d_\alpha$ 个 D6-膜表示，它们处于固定的 $x^{4,5,6}$ 位置，且 $x^6$ 的值落在区间 $I_\alpha$ 内。我们用 $K$ 记诸 $k_\alpha$ 的最大值。当我们把 D6-膜包括进来时，左右两端的半无限 D4-膜并不以任何有趣的方式改变分析，所以本节余下部分我们将把它们略去。与之前一样，对 $\alpha =1,\dots, n$，$k_\alpha$ 个 D4-膜沿区间 $I_\alpha$ 张在 NS5-膜之间。我们假设最初没有 D4-膜终止于 D6-膜上。在弱 IIA 耦合极限下，D4-D6 弦提供基本物质。当这 $k_\alpha$ 个 D4-膜重合时，这些基本超多重态的质量参数由 D4-膜与 D6-膜的 $v$ 坐标之差给出。

我们刚刚描述的位形，其 M 理论提升是这样得到的：先把 D6-膜的位形提升为多中心 Taub-NUT 几何 $Q$，然后在库仑分支的真空上把 D4-膜与 NS5-膜提升为单独一张 M5-膜，其世界体为 $\IR^{1,3} \times \tilde \Sigma \times \{\CP\}$，其中 $\tilde \Sigma \subset Q$ 是一张全纯曲面。本节的目标是再一次与柱面上的 Hitchin 系统取得联系，并说明基本物质如何改变 Higgs 场上的边界条件。

结果表明，做到这一点有若干种不同的方式，因为同一膜系统可以经受某些形变，而这些形变给出相同的红外四维场论、却给出不同的中间六维构造。详细分析稍嫌繁复，所以我们先预揭主要结果。

对线性共形箭图存在一个典范选择，它仍然给出柱面上带正则奇点的 $SU(K)$ Hitchin 系统。$t=0$ 与 $t=\infty$ 处的边界条件，即 \eqref{eq:SimpleAtZero} 与 \eqref{eq:SimpleAtInfty}，有轻微的修改，因为 $v^{(0)}_i$ 与 $v^{(n+1)}_j$ 不再是通用的了。取而代之，$K$ 个 $v^{(0)}_i$ 被分成若干相同值的块，使得每一块由 $\alpha$ 个相同的 $v^{(0)}_i$ 组成、对应箭图第 $\alpha$ 个节点处的一个基本味；而每一块由 $\alpha$ 个相同的 $v^{(n+1)}_i$ 组成、对应箭图第 $(n+1-\alpha)$ 个节点处的一个基本味。例如 $K=3$ 时，我们有若干 $SU(3)$ 规范节点的箭图，可能在一端或两端各有一个 $SU(2)$ 规范群。若取 $k_1=2$、$k_2=k_3=3$，则 $d_1=d_2=1$、$d_3=3$，于是第一个节点处有一个基本场，第二个节点处又有一个。这时 $t=\infty$ 处正则奇点的留数将有两个相同的本征值。

这个典范 Hitchin 系统的特征是：四维理论的 $U(1)$ R 对称性与旋转 $v$ 的 $U(1)$ 对称性等同。其他的选择导致柱面上（秩高于 $K$ 的）Hitchin 系统，在 $t=0,\infty$ 之外没有穿刺点，而在 $t=0$ 或 $t=\infty$ 处有野分歧。四维理论的 $U(1)$ R 对称性被等同于旋转 $v$ 与旋转 $t$ 的两个 $U(1)$ 对称性的某个组合。同一四维场论——从而同一超凯勒模空间——竟由不同的 Hitchin 系统描述，这是令人惊讶的。关于物质我们将在第 3.2.8 节% TODO 合稿改 \ref{sec:surprise}
进一步讨论。

先来回忆 $Q$ 上的度量。$Q$ 上的度量由位置 $\vec r_a = (r_{1,a}, r_{2,a}, r_{3,a}) \in \IR^3$ 决定，其中 $a$ 跑遍某个指标集。若把 $Q$ 描述为 $\IR^3$ 上的圆丛，则度量为
\begin{equation}\label{eq:ch05-dsTN}
ds^2_{TN} = V^{-1} (d\psi +A)^2 + V d\vec r^2
\end{equation}
其中 $\psi \sim \psi + 4\pi$，
\begin{equation}\label{eq:ch05-V}
V =1 + \sum_{a} \frac{1}{\vert \vec r - \vec r_a\vert}
\end{equation}
且 $dA = *dV$。完整的 M 理论度量为
\begin{equation}\label{eq:ch05-dsM}
ds^2 = dx^\mu dx_\mu + \frac{R_{11}^2}{4} ds^2_{TN} + dx^2_{7,8,9}
\end{equation}
我们等同 $s = x^6 + i x^{10} = \half ( r_3 + i \psi)$ 与 $r_1+i r_2 = 2 \frac{\ell}{R_{11} } v$。Taub-NUT 空间 $Q$ 带有超凯勒结构，但其中有一个特殊的复结构，使得坐标 $v$ 是全纯的。在这个复结构下，流形 $Q$ 作为复流形有方程
\begin{equation}\label{eq:MultiCenterTNa}
UW  = \prod_a (v-v_a)
\end{equation}
参数 $v_a$ 是复结构参数，而 $x^6_a$ 是 $Q$ 上超凯勒度量的凯勒参数。

现在转向复流形 \eqref{eq:MultiCenterTNa} 中全纯曲线 $\tilde \Sigma$ 的描述。这条曲线可以描述为一个多项式方程 $\tilde F(W,v)=0$。或者在另一个坐标卡中，利用 $U=W^{-1} \prod_a (v-v_a)$，该方程也可以改写为 $\tilde G(U,v)=0$，其中 $\tilde G$ 是多项式。文献 \cite{Witten:1997sc} 表明，「左右两端都没有半无限 D4-膜」这一约束导致 $\tilde F(W,v)$ 具有如下结构：对每个区间 $I_\alpha := \{ x^6: {\rm Re}(s_{\alpha -1}) < x^6 < {\rm Re}(s_\alpha)\}$，联系一个次数为 $d_\alpha$ 的多项式：
\begin{equation}\label{eq:ch05-Jalpha}
J_\alpha(v):= \prod_{x^6_a\in I_\alpha} (v-v_a), \qquad\qquad \alpha=1,\dots, n.
\end{equation}
现在引入
\begin{equation}\label{eq:Aalphdef}
A_0(v) = c_0
\end{equation}
\begin{equation}\label{eq:Aalphdef-2}
A_\alpha(v) = c_\alpha g_\alpha(v) \prod_{\beta =1}^{\alpha }
J_\beta^{\alpha-\beta}(v), \qquad\qquad \alpha = 1, \dots, n
\end{equation}
\begin{equation}\label{eq:Aalphdef-3}
A_{n+1}(v)  = c_{n+1}   \prod_{\beta =1}^{n}
J_\beta^{n+1-\beta}(v),
\end{equation}
其中 $c_\alpha$（$\alpha = 0, \dots, n+1$）是非零常数，而 $g_{\alpha}(v)$ 是 $v$ 的次数为 $k_\alpha$ 的\emph{首一}多项式，$\alpha = 1, \dots, n$。（同时约定 $g_0 = g_{n+1}=1$ 也方便。）于是我们有曲线（文献 \cite{Witten:1997sc}，式 (3.23)）
\begin{equation}\label{eq:TNCurveii}
\tilde F(W,v) =  \sum_{\alpha=0}^{n+1} A_\alpha(v) W^{n+1-\alpha}=0
\end{equation}
或等价地
\begin{equation}\label{eq:TNCurveiii}
\tilde G(U,v) =  c_{n+1}U^{n+1} + c_n g_n(v) U^n + \sum_{\alpha=0}^{n-1}
g_{\alpha}(v)\left( \prod_{\beta = \alpha+1}^n J_\beta(v)^{\beta-\alpha}\right) U^\alpha =0
\end{equation}
此外，$g_{\alpha}(v)$ 的系数在文献 \cite{Witten:1997sc} 中被解释为 $SU(k_\alpha)$ 规范群通常的序参量 $\langle {\Tr}(\Phi^{(\alpha)})^s\rangle$。常数 $c_\alpha$ 当然可以乘上同一个非零因子重新标度，因此应视为射影空间中的一个点。它们编码规范耦合。（在弱耦合区域，$c_\alpha$ 可以用 $t_\alpha$ 的第 $\alpha$ 个初等对称函数表出。）

乍看之下，Taub-NUT 构造似乎不会导致 Hitchin 系统有用的情形，因为它与我们先前那幅「$T^*C$ 上 $K$ 张五膜缠绕 $C$、并带横向缺陷」的图景颇为不同。然而，正如我们现在要解释的，我们可以取 Taub-NUT 几何的一个极限，在其中重新得到 Hitchin 描述。

曲线 $\tilde \Sigma$ 将再次被解释为 SW 曲线。回忆多项式 $\tilde F(W,v)$ 与凯勒参数 $x^6_a$ 无关，$Q$ 的复结构亦然。正如我们在第 3.1.5 节% TODO 合稿改 \ref{subsubsec:KineticTerms}
以及在式 (3.54)% TODO 合稿改 \eqref{eq:KEreduxii}
及其后各式中看到的，动能项与 SW 微分只依赖于复结构，因而不受 $x^6_a$ 改变的影响：两个仅相差 D6-膜沿 $x^6$ 方向平移的理论，由同一个红外不动点描述。另一方面，在把 D6-膜远远移到 NS5-膜系统左侧或右侧的极限下，该位形与不带 D6-膜的位形极为相似，因为当 $\vec r$ 取有限值时，$\vert x^6_a\vert \to \infty$ 显然给出 $V\to 1$。

现在把 D6-膜分成两个不相交的子集 $\CL \amalg \CR$，并考虑一个极限：$a \in \CR$ 的 D6-膜远远右移，而 $a\in \CL$ 的 D6-膜远远左移。在附录 D% TODO 合稿改 \ref{app:HoloTN}
中我们证明，适当归一化之后 $W$ 的极限值是
\begin{equation}\label{eq:W-t-limit}
W \to t \prod_{a\in\CL}(v-v_a).
\end{equation}
把它代入 \eqref{eq:TNCurveii} 并化简。为此注意每个多项式 $J_\alpha(v)$ 都按照移向左方与移向右方的 D6-膜分解：$J_\alpha(v) = J_{\alpha,L}(v) J_{\alpha,R}(v)$。在提出因子 $\prod_{\alpha} J_{\alpha,L}^{n+1-\alpha}$ 之后，我们剩下
\begin{equation}\label{eq:Fhatcurve}
\hat F(t,v) = \sum_{\alpha=0}^{n+1} \hat g_{\alpha}(v) t^{n+1-\alpha},
\end{equation}
其中
\begin{equation}\label{eq:hatgee}
\hat g_0(v)  =  c_0 \prod_{\beta=1}^n J_{\beta,L}^\beta
\end{equation}
\begin{equation}\label{eq:hatgee-2}
\hat g_{\alpha}(v)  =  c_{\alpha} g_{\alpha}(v) \prod_{\beta=1}^{\alpha} J_{
\beta,R}^{\alpha-\beta} \prod_{\beta=\alpha+1}^{n}
J_{\beta,L}^{\beta-\alpha}\qquad\qquad \alpha =1, \dots, n-1
\end{equation}
\begin{equation}\label{eq:hatgee-3}
\hat g_n(v)  =  c_{n} g_{n}(v) \prod_{\beta=1}^{n} J_{
\beta,R}^{n-\beta}
\end{equation}
\begin{equation}\label{eq:hatgee-4}
\hat g_{n+1}(v) =  c_{n+1} \prod_{\beta=1}^{n}
J_{\beta,R}^{n+1-\beta}.
\end{equation}

对 $\hat g_{\alpha}(v)$ 中的各因子，有一个基于 Hanany–Witten 效应 \cite{Hanany:1996ie} 的漂亮解释。在我们的语境中，Hanany–Witten 效应说的是：对每个 D6-膜，「链接数」
\begin{equation}\label{eq:ch05-linking}
(L-R) - \half (l -r )
\end{equation}
是常数。这里 $l$（$r$）是该 D6-膜左侧（右侧）NS5-膜的数目，而 $L$（$R$）是从左（右）侧终止于该 D6-膜上的 D4-膜数目。特别地，若一个 D6-膜在固定 $v$ 下沿 $x^6$ 方向穿过一个 NS5-膜，则张在二者之间的一个 D4-膜被产生出来。在我们把 $a\in \CR$ 的 D6-膜远远右移、把 $a\in \CL$ 的 D6-膜远远左移之后，每个区间都由「自由」D4-膜与「冻结」D4-膜组成。\footnote{「冻结」这个形容词指的是文献 \cite{Hanany:1996ie} 中讨论的所谓 $S$-规则。由 D6-膜的横向运动产生、位于区间 $I_\alpha$ 中的 D4-膜，在保持超对称的同时不能在 $v$ 方向上独立移动。} $\hat g_\alpha$ 中的因子 $g_{\alpha}(v)$ 解释「自由」D4-膜。$g_{\alpha}(v)$ 的系数按通常方式对应可归一的形变。因子 $\prod_{\beta=1}^{\alpha} J_{R,\beta}^{\alpha-\beta}$ 解释区间 $I_\alpha$ 中因 D6-膜向右运动而产生的「冻结」D4-膜，而因子 $\prod_{\beta=\alpha+1}^{n} J_{\beta,L}^{\beta-\alpha}$ 解释区间 $I_\alpha$ 中因 D6-膜向左运动而产生的「冻结」D4-膜。

计算多项式 $\hat g_{\alpha}$ 的次数 $\hat k_\alpha$ 有一定意义。为此注意：一个 D6-膜以任一方向穿过区间 $I_\alpha$ 之后，$-2k_\alpha + k_{\alpha -1} + k_{\alpha +1} + d_\alpha$ 的净值保持不变。因此，既然我们初始的位形把这些值全取为零（为了共形性），就可推出：在所有 D6-膜都被移到最左或最右之后，有 $-2\hat k_\alpha + \hat k_{\alpha+1} + \hat k_{\alpha-1}=0$，$1\leq \alpha \leq n$。这里 $\hat k_0$ 是终止于左侧 D6-膜上的 D4-膜数目，即 $\prod J_{\beta,L}^\beta$ 的次数；类似地，$\hat k_{n+1}$ 是终止于右侧 D6-膜上的 D4-膜数目，即 $\prod J_{\beta,R}^{n+1-\beta}$ 的次数。利用附录 E% TODO 合稿改 \ref{app:ConvexKay}
的结果我们看到 $\frac{(\hat k_{n+1}  - \hat k_0 )}{n+1}=r$ 是非负整数，且
\begin{equation}\label{eq:solvekhat}
\hat k_\alpha = \hat k_0 + \alpha r.
\end{equation}

把 D6-膜向左、向右移动的方式有很多。一个极端是把它们全移到右边，于是 $\hat k_{0}=0$，而 $\hat k_\alpha$ 作为 $\alpha$ 的函数线性增大。当然，另一个极端是把所有 D6-膜放在左边，此时 $\hat k_\alpha$ 作为 $\alpha$ 的函数线性减小。在左右移动 D6-膜的种种方式之中，有一个典范选择，我们将称之为「平衡情形」。为描述它，我们再次求助于附录 E% TODO 合稿改 \ref{app:ConvexKay}
：注意 $k_\alpha$ 作为 $\alpha$ 的函数增大，直到某个最大值 $K$ 在某个 $\alpha=\alpha_{-}$ 处达到，然后 $k_\alpha=K$ 一直保持不变，直到某个 $\alpha=\alpha_+\geq \alpha_-$，此后 $k_\alpha$ 对 $\alpha > \alpha_+$ 单调减小。这提示 D6-膜的一种典范移动：把区间 $\alpha < \alpha_-$ 中的所有 D6-膜带到左边，把区间 $\alpha > \alpha_+$ 中的所有 D6-膜带到右边。在区间 $I_{\alpha_-}$ 中我们把 $k_{\alpha_-}-k_{\alpha_--1}$ 个膜带到左边，在区间 $I_{\alpha_+}$ 中我们把 $k_{\alpha_+}-k_{\alpha_++1}$ 个膜带到右边。不难证明这导致 $\hat k_{\alpha_--1} = \hat k_{\alpha_-} = \cdots = \hat k_{\alpha_+} = \hat k_{\alpha_++1}$。既然 $\hat k_\alpha$ 线性增长，就推出 $r=0$，从而全部 $\hat k_\alpha$ 都等于某个公共值 $\hat K$。

这种 D6-膜的「平衡」移动把系统映到一个与上一小节极为相似的系统。D6-膜左右分布的其余可能情形我们留到下一小节讨论，而本节余下部分只关注这个典范选择。此时曲线 \eqref{eq:Fhatcurve} 是 (3.46)% TODO 合稿改 \eqref{eqn:sw1}
的特例，其中 $q_\alpha(v)$ 形如 $\hat g_\alpha(v)$。与之前一样，我们可以把 $\hat F(t,v)$ 整理成 $v$ 的 $\hat K$ 次多项式，其形式为 (3.45)% TODO 合稿改 \eqref{eq:swcurve1}
：
\begin{equation}\label{eq:ch05-Fhatmonic}
\hat F(t,v) = v^{\hat K} \prod_{\alpha=0}^n(t-t_\alpha) + \sum_{i=1}^{\hat
K} v^{\hat K-i} p_i(t) 
\end{equation}
除以 $\prod_{\alpha=0}^n(t-t_\alpha)$，我们得到 $v$ 的首一多项式，其系数 $R_i(t) =
\frac{p_i(t)}{\prod_{\alpha=0}^n(t-t_\alpha)}$ 是 $t$ 的有理函数。现在把它解释为一条谱曲线的方程
\begin{equation}\label{eq:ch05-detspectral}
\det(v- t \varphi_t) =  v^{\hat K}   + \sum_{i=1}^{\hat K} v^{\hat
K-i} R_i(t)
\end{equation}
其中 Higgs 场 $\varphi_t$ 取值于 $u(\hat K)$。我们的目标现在是理解系数 $\hat g_\alpha(v)$ 的这种特殊结构，对 Higgs 场 $\varphi_t$ 在缺陷处的边界条件意味着什么。

于是，让我们分析 \eqref{eq:Fhatcurve} 的 $\hat K$ 个根 $v_i(t)$ 的行为。既然 $t$ 作为 $v$ 的函数有 $n+1$ 个根 $t(v)$，而 $v$ 可以自由趋于无穷，就必然存在 $(n+1)$ 个值 $t_\alpha$ 使 $v\to \infty$。一般地，$v$ 在 $t_\alpha$ 处有一个单极点。这些单极点意味着（在移走质心之后）$\varphi_t$ 有一个一阶极点，其留数恰如 \eqref{eq:SimplePoleResidue} 所述。这正是我们先前遇到的情形。

另一方面，$t\to 0$ 或 $t\to \infty$ 时根的行为，需要我们此前分析的某种精细化。$t\to 0$ 时根趋于 $\hat g_{n+1}(v)$ 的根，而 $t\to \infty$ 时趋于 $\hat g_0(v)$ 的根。这些根都在有限的 $v$ 值处。然而，$\hat g_0$ 中 $\beta>1$ 的因子 $J_{\beta,L}$（以及 $\hat g_{n+1}$ 中 $\beta < n+1$ 的因子 $J_{\beta,R}$）导致重根。一般地，系数 $q_0(v)$（或 $q_{n+1}(v)$）在 (3.46)% TODO 合稿改 \eqref{eqn:sw1}
中存在重根，意味着若干 D4-膜终止于同一个 D6-膜上。这里有一个重要的区别：一类重根来自把若干附着于单个 D4-膜的 D6-膜移到同一个 $v$ 坐标，另一类重根则来自 Hanany–Witten 效应。在前一情形，$q_0(v)$ 的根一般不同于 $\alpha>1$ 时 $q_\alpha(v)$ 的根。在后一情形，(3.90)% TODO 合稿改 \eqref{eq:hatgee}
及其后各式的结构表明，来自因子 $J_{\beta,L}$ 的重根同时也是 $\hat g_{\alpha}$（$\alpha < \beta$）的（重）根。于是，一般地，当 $t \to \infty$ 时，若 $q_0(v)$ 有一个 $\beta$ 重根 $v_*$，而 $v_*$ 不是 $q_1(v)$ 的根，则 $\beta$ 个根 $v_i(t)$ 的行为是
\begin{equation}\label{eq:FracRootLimit}
v_i(t) \to v_* + \frac{\xi_i}{t^{1/\beta}} + \cdots
\end{equation}
（当 $t\to\infty$），其中 $\xi_i$ 是与 $t,v$ 无关的常数。而另一方面，对来自 (3.90)% TODO 合稿改 \eqref{eq:hatgee}
中因子 $J_{\beta,L}$ 的重根，相应的 $\beta$ 个根 $v_i(t)$ 的行为是
\begin{equation}\label{eq:WholeRootLimit}
v_i(t) \to v_* + \frac{\xi_i}{t}+ \CO(t^{-2})
\end{equation}
对右侧 D6-膜所联系的根在 $t\to 0$ 时的行为，有完全类似的陈述。

我们刚刚描述的 $t\to \infty$ 时根 $v(t)$ 的行为，对 $t\to\infty$ 时 Higgs 场 $\varphi_t dt$ 的边界条件有含义。这个 Higgs 场的行为将是
\begin{equation}\label{eq:ch05-phiAsym}
\varphi_t \to \frac{R}{t} + \frac{R_2}{t^2} + \cdots
\end{equation}
把 $\hat F(t,v) = \hat g_0(v) t^{n+1} + \cdots $ 与 $\prod_0^n (t-t_\alpha) \det(v-t\varphi_t)$ 等同起来，并取 $t\to \infty$，我们看到留数 $R = t \varphi_t\vert_{t=\infty}$ 的特征多项式恰是 $\det(v-R) = \hat g_0(v)$。既然它有重根，我们就必须考虑 $R$ 具有非平凡 Jordan 型的可能。

现在我们断言：非平凡的 Jordan 型导致根的行为如 \eqref{eq:FracRootLimit}，而若 $R$ 半单，则根的行为如 \eqref{eq:WholeRootLimit}。我们可以如下证明。为简单起见，设特征多项式 $\det(v-R) = v^\beta$，其中整数 $\beta > 1$。现在考虑微扰 $R(t) = R + C/t $，其中 $t$ 很大而 $C$ 是\emph{通用}的。若 $R$ 半单，则 $R=0$，而 $R(t)$ 的本征值就是 $c_i/t$，其中 $c_i$ 是 $C$ 的互异本征值。另一个极端，设 $R$ 是一个大小为 $\beta$ 的 Jordan 块，我们把它记作 $N_\beta$。\footnote{即用矩阵单位来表示，$N_\beta = e_{1,2} + e_{2,3} + \cdots + e_{\beta-1,\beta}$。} 通过 $g= \exp[\epsilon/t]$ 作规范变换，我们可以把 $R(t)$ 化为
\begin{equation}\label{eq:ch05-Jordanpert}
R(t) \to N_\beta + \frac{1}{t}\left( \sum_{i=1}^\beta \mu_i
e_{\beta,i} \right) + \CO(1/t^2)
\end{equation}
其中 $\mu_i$ 为某些常数。这里 $e_{i,j}$ 是矩阵单位，即第 $i$ 行第 $j$ 列的元素为 $1$、其余元素为零的矩阵。另一方面，
\begin{equation}\label{eq:UsefulDet}
\det\left[ v - N_\beta -\frac{1}{t}\left( \sum_{i=1}^\beta \mu_i
e_{\beta,i} \right)\right]= v^{\beta} +
\frac{\mu_{\beta}}{t}v^{\beta-1} + \frac{\mu_{\beta-1}}{t}
v^{\beta-2} + \cdots + \frac{\mu_1}{t}
\end{equation}
由初等对称函数与幂和函数之间的标准关系推出，根满足 $v_i
\sim c_i t^{-1/\beta} + \CO(t^{-2/\beta}) $。在中间情形，即 $R$ 有若干个 Jordan 块时，根将以 $t$ 的分数幂衰减，而该分数幂由 $R$ 的最大 Jordan 块控制。

结论是：$t\to \infty$ 处 Higgs 场的边界条件呈块对角形式：
\begin{equation}\label{eq:RegSingBalanced}
\varphi_t dt \to  \frac{dt}{t} {\rm Diag}\{ \cdots  v^{(\beta)}_1
1_{\beta},\dots,v^{(\beta)}_{d_{\beta,L}} 1_{\beta}, \cdots \} +
\CO(t^{-2})
\end{equation}
其中 $1_{\beta}$ 是对角矩阵，而 $v^{(\beta)}_1
,\dots,v^{(\beta)}_{d_{\beta,L}}$ 是 $J_{\beta,L}$ 的根（为简单起见假定互异）。更简单地说，这些边界条件保持 $U(\hat K)$ 规范群的一个子群
\begin{equation}\label{eq:ch05-subgroup}
\prod_{\beta=1}^n (U(\beta))^{d_{\beta,L}}
\end{equation}
对 $t\to 0$ 的极限以及右侧的膜群，有完全平行的论述。

总之，酉群的共形线性箭图给出 $\IC \IP^1$ 上的 Hitchin 系统，它有两个一般类型的正则奇点，分别由 $K$ 的两个分拆 $K = \sum \beta d_{\beta,L}$ 与 $K = \sum \beta d_{\beta,R}$ 标记，此外还有任意多个「基本」奇点，联系于分拆 $K = (1) + (K-1)$。

\subsubsection{弯膜与不规则奇点}\label{subsubsec:Brane-bending}

我们继续研究上一小节的共形不变线性箭图，但现在考虑以「非平衡」方式移动 D6-膜、使得 $\hat k_{n+1} - \hat k_0$ 非零的情形。

这时根的行为极为不同。回忆 $\hat g_{\alpha}$ 的次数是 $\hat k_\alpha = \hat k_0 + \alpha r$。为简单起见设 $r>0$，于是按 \eqref{eq:solvekhat}，$\hat g_{n+1}= c_{n+1} v^{\hat k_{n+1}} + \cdots $ 含有 $v$ 的最高次幂。那么，既然 $v\to \infty$ 时 $\hat g_{n+1}$ 支配其他所有系数 $\hat g_\alpha(v)$，就推出没有任何根 $v(t)$ 在有限的 $t$ 值处趋于无穷。

只有两个 $t$ 方向的奇点这一事实，也可以通过理论的 $U(1)_R$ 对称性更物理地看出来。这个对称性旋转 $v \to e^{i \theta} v$，从而改变质量 $v_a \to e^{i \theta} v_a$。然而它不能改变耦合常数，因此 $c_\alpha \to c_\alpha$。在平衡情形，曲线 $\hat F(t,v)=0$ 在固定 $t$ 下对这个标度变换不变。另一方面，在非平衡情形这不再成立，我们必须同时标度 $(t,v) \to (e^{i r\theta} t, e^{i \theta} v)$。于是，位于有限点 $t_\alpha\in \IC^\times$ 的奇点将与方程不相容。

其次，考虑 $t\to 0$ 时的根。这里 $v_i(t)$ 简单地趋近 $\hat g_{n+1}(v)$ 的 $\hat k_{n+1}$ 个根；而如我们已经讨论过的，它们的行为是 $v_i(t) = v_i + \xi_{i,s} t + \cdots$。

另一方面，$t \to \infty$ 时根的行为更复杂。有 $\hat k_0$ 个根的行为是 $v_i(t) \to v_i + \frac{\xi_{i,a}}{t}+ \CO(1/t^2)$，其中 $v_i$ 是 $\hat g_0(v)$ 的根。由于 $\hat k_{n+1} > \hat k_0$，这并不能解释全部根。此外还有 $(n+1)r$ 个根使 $t, v$ 同时趋于无穷。此时保留 $\hat F(t,v)$ 中的首阶项，我们看到这些 $(t,v)$ 同时趋于无穷的根必定渐近于
\begin{equation}\label{eq:BigRoots}
0 = c_{n+1} v^{(n+1)r} + c_n t v^{nr} + \cdots + c_1 t^n v^r + c_{0}
t^{n+1} := c_{n+1} \prod_{i=1}^{n+1} (v^r - \nu_i t).
\end{equation}
的根。

\insfig{3.3in}{brane-bending}{共形但非平衡情形（$\hat k_{n+1} - \hat k_0 = (n+1) r > 0$）下的弯膜。注意有 $n+1$ 束、每束 $r$ 张膜当 $t\to \infty$ 时趋于无穷。这里 $r=3$，$n+1=4$。在 $t\to 0$ 处有 $\hat k_{n+1}=17$ 张水平膜。}

注意我们现在发现了「弯膜」现象：某些根 $v(t)$ 当 $t\to \infty$ 时趋于无穷。为把这一情形画出来，注意 $\Sigma$ 在 $(-\log\vert t\vert, \log \vert v \vert)$ 平面上的投影渐近于图 \ref{fig:brane-bending} 所示的形状。

现在考虑映射到 Hitchin 系统。此时我们如下认定 Higgs 场的特征多项式
\begin{equation}\label{eq:ch05-detchar}
\det(v - t \varphi_t) = \frac{1}{c_{n+1}} \left( \hat g_{n+1}(v) + t
\hat g_n(v) + \cdots + t^{n+1} \hat g_0(v) \right)
\end{equation}
其中 $\varphi_t \in u(\hat k_{n+1})$。由上面关于根的讨论，Higgs 场的全部奇点必定位于 $t=0$ 或 $t=\infty$。我们此前的讨论适用于 $t\to 0$ 处的根，因此 $\varphi_t$ 在 $t \to 0$ 处有一个正则奇点，如 \eqref{eq:RegSingBalanced} 所述。

相反，对 $t\to \infty $ 我们发现 $\varphi_t$ 呈块对角形式。有一个 $\hat k_0 \times \hat k_0$ 的块对应一个正则奇点，其留数由 $\hat g_0(v)$ 的根给出。我们把 \eqref{eq:BigRoots} 解释为：此外还有 $(n+1)$ 个块，以 $i=1,\dots, n+1$ 标记，其极限行为是
\begin{equation}\label{eq:ch05-blocklimit}
\det(v-t\varphi_t) \sim v^r - \nu_i t.
\end{equation}
（矩阵被理解为限制在第 $i$ 个块上。为避免记号累赘，我们不把它标出。）于是，在非平衡情形，$t\to \infty$ 处的边界条件对应一个\emph{不规则奇点}，其中 $t\varphi_t$ 的 $(n+1)r$ 个本征值的行为是
\begin{equation}\label{eq:ch05-eiginfty}
\omega^s (\nu_i t)^{\frac{1}{r}} (1 + \CO(1/t)) \qquad\qquad
s=1,\dots, r;\qquad  i = 1, \dots, n+1,
\end{equation}
其中 $\omega$ 是 $r$ 次本原单位根。

我们也可以考虑相反的情形 $\hat k_0 > \hat k_{n+1}$。此时我们反转 $r$ 的符号并取 $r = (\hat k_0-\hat k_{n+1})/(n+1)>0$。这时 $t \varphi_t$ 在 $t\to \infty$ 处有正则奇点，而在 $t \to 0$ 处有不规则奇点。这个不规则奇点涉及 $(n+1)$ 个块，每块的本征值为
\begin{equation}\label{eq:ch05-eigzero}
\omega^s (\frac{\nu_i}{t})^{\frac{1}{r}} (1 + \CO(t)) \qquad\qquad
s=1,\dots, r;\qquad  i = 1, \dots, n+1
\end{equation}
其中 $\omega$ 是 $r$ 次本原单位根。

现在讨论这些条件对 Higgs 场 $\varphi_t$ 在 $t\to 0,\infty$ 处行为的含义。我们集中看 $t\to 0$，并取 $r = (\hat k_0-\hat k_{n+1})/(n+1)>0$。利用 \eqref{eq:UsefulDet} 我们看到，通过一个\emph{复}规范变换，可以把 $\varphi_t$ 化为
\begin{equation} \label{eq:phi-formal}
\varphi_t \to \frac{1}{t} \left( \frac{\nu_i}{t} e_{r1} + N_r +
\CO(t) \right) \qquad t \to 0 \end{equation}
然而，要正确定义 Hitchin 系统，我们必须给出 $t\to 0$ 处 Hitchin 方程的一个实际解。这可以逐块进行。对对应不规则奇点的那些块，在对角规范下容易给出这样的解，
\begin{equation}\label{eq:ch05-diagformal}
\varphi_t = \frac{1}{t} \left(\frac{\nu_i}{t}\right)^{1/r} {\rm
Diag} \{ \omega, \omega^2, \dots, \omega^r \}
\end{equation}
问题在于这不是单值的。我们可以通过一个奇异规范变换使 $\varphi_t$ 单值，
\begin{equation}\label{eq:ch05-gaugeeta}
g = \eta S \eta^{-1}
\end{equation}
其中 $S_{ab} = \frac{1}{\sqrt{r}} e^{\frac{2\pi i }{r} ab}$，而 $\eta$ 由下式给出
\begin{equation}\label{eq:ch05-eta}
\eta = \left(\frac{t}{\bar t} \right)^{\frac{r-1}{4r}} {\rm Diag}\{
1,\left(\frac{\bar t}{ t} \right)^{\frac{1}{2r}},\left(\frac{\bar
t}{ t} \right)^{\frac{2}{2r}},\dots, \left(\frac{\bar t}{ t}
\right)^{\frac{r-1}{2r}} \}
\end{equation}
为使规范场保持单值，我们取
\begin{equation}\label{eq:ch05-Aeta}
A = - d\eta \eta^{-1}
\end{equation}
（在对角规范下）。结论是：我们的 Hitchin 方程的单值渐近解（对第 $i$ 个块）是
\begin{equation}\label{eq:ISPatInfty}
\varphi_t \to \frac{\nu_i^{1/r}}{t} \frac{1}{\vert t\vert^{1/r}}
\left( \frac{\bar t}{\vert t\vert } e_{r1} + N_r + \CO(t) \right)
\end{equation}
并伴随
\begin{equation}\label{eq:GaugeFldInfty}
A  =  \left( \sum_{j=1}^r \frac{2(j-1)-(r-1)}{4r} e_{jj} \right)
\left( \frac{dt}{t} - \frac{d\bar t}{\bar t}\right).
\end{equation}

\subsubsection{渐近自由情形}

最后我们转向带基本物质的一般渐近自由箭图。对这种箭图，$\beta$ 函数系数 $b_{\alpha} \leq 0$（$1\leq \alpha \leq n$），且至少有一个 $\alpha$ 使 $b_\alpha<0$。物理上，这类箭图规范理论可以通过把共形线性箭图中某些基本场的质量取到无穷、使这些基本场脱耦而得到。它们脱耦之后，留下一个渐近自由理论。

脱耦程序可以在 SW 曲线的层次上优雅地完成。在每个区间 $I_\alpha$ 中我们选一组中心 $v_a$（$a\in S_\alpha$），把它们送到无穷。然后我们以使得 $A_\alpha(v)$ 在式 \eqref{eq:Aalphdef}–\eqref{eq:Aalphdef-3} 中保持有限的方式令 $c_\alpha \to 0$。不过，既然这里我们最关心对 Hitchin 系统的应用，我们考虑的是在把 D6-膜移到大的 $\vert x_a^6\vert$ \emph{之后}的脱耦极限。因此我们按照待脱耦的 D6-膜先被移到最左还是最右，把集合分成 $S_\alpha = S_{\alpha,L} \amalg S_{\alpha,R}$。在这个极限下我们有
\begin{equation}\label{eq:ch05-JLlimit}
J_{\alpha,L}(v) \to \check J_{\alpha,L}(v) \prod_{a\in S_{\alpha,L}} (-v_a) (1 +
\CO(v/v_a)),
\end{equation}
其中 $\check J_{\alpha,L}(v)$ 是未被送到无穷的 D6-膜的因子，$J_{\alpha,R}$ 类似。于是我们令 $c_\alpha$ 趋于零，使得
\begin{equation}\label{eq:ch05-cctilde}
\check c_{\alpha} = c_\alpha \prod_{\beta=1}^\alpha \left( \prod_{a\in
S_{\beta,R}}(-v_a)\right)^{\alpha -\beta}\prod_{\beta=\alpha+1}^n \left( \prod_{a\in
S_{\beta,L}}(-v_a)\right)^{\beta -\alpha}
\end{equation}
保持有限。现在我们有方程 \eqref{eq:Fhatcurve} 与 \eqref{eq:hatgee} 的类似物：
\begin{equation}\label{eq:Fhatcurvechck}
\check F(t,v) = \sum_{\alpha=0}^{n+1} \check g_{\alpha}(v) t^{n+1-\alpha}
\end{equation}
其中
\begin{equation}\label{eq:hatgeechck}
\check g_0(v) = \check c_0 \prod_{\beta=1}^n \check J_{\beta,L}^\beta
\end{equation}
\begin{equation}\label{eq:hatgeechck-2}
\check g_{\alpha}(v) = \check c_{\alpha} g_{\alpha}(v) \prod_{\beta=1}^{\alpha} \check J_{
\beta,R}^{\alpha-\beta} \prod_{\beta=\alpha+1}^{n} \check
J_{\beta,L}^{\beta-\alpha}\qquad\qquad \alpha =1, \dots, n-1,
\end{equation}
\begin{equation}\label{eq:hatgeechck-3}
\check g_{n}(v) = \check c_{\alpha} g_{\alpha}(v) \prod_{\beta=1}^{n} \check J_{
\beta,R}^{\alpha-\beta},
\end{equation}
\begin{equation}\label{eq:hatgeechck-4}
\check g_{n+1}(v)  =  \check c_{n+1} \prod_{\beta=1}^{n} \check
J_{R,\beta}^{n+1-\beta}.
\end{equation}
与前一情形的唯一差别是：现在 $\check b_\alpha = -2\check k_\alpha + \check k_{\alpha + 1} + \check k_{\alpha -1}$ 对一切 $1\leq \alpha \leq n$ 都 $\leq 0$，而且可以严格小于零。这里 $\check k_\alpha$ 是 $\check g_{\alpha}$ 的次数。

再次援引附录 E% TODO 合稿改 \ref{app:ConvexKay}
，点 $(\alpha, \check k_\alpha)$ 定义一条凸折线。$\check k_\alpha$ 在 $0 \leq \alpha \leq \alpha_-$ 上严格递增。最大值 $\check K := \check k_{\alpha_-}$ 在 $\alpha_- \leq \alpha \leq \alpha_+$ 上取到。此后 $\check k_\alpha$ 在 $\alpha \geq \alpha_+$ 上严格递减。若 $\check k_1 - \check k_0 =0$，则 $\alpha_-=0$；若 $\check k_1 - \check k_0 < 0$，则 $\alpha_+=0$。为简单起见，下面我们只描述 $\check k_1 - \check k_0 > 0$ 的情形。

\insfig{4.7in}{brane-bending-2}{$\check F(t,v)$ 的牛顿多边形的边界支配 M5-膜曲线的渐近行为。这里把它投影到 $(-\log\vert t \vert, \log \vert v \vert )$ 平面上。}

现在可以把这条折线看作 $\check F(t,v)$ 的牛顿多面体的边界。据此我们可以提取方程 $\check F(t,v)=0$ 之根的渐近行为。做法是多面体边界面上各项的平衡。为简单起见，假定一个面由单一线段组成，我们有：
\begin{equation}\label{eq:ch05-balance}
\check c_\alpha t^{n+1-\alpha} v^{\check k_\alpha} + \check c_{\alpha +1} t^{n-\alpha} v^{\check k_{\alpha +1} }
\cong 0
\end{equation}
对 $0 \leq \alpha \leq \alpha_--1$ 或 $\alpha_+\leq \alpha \leq n$。每个 $\alpha$ 值给出 $\vert \check k_{\alpha+1} - \check k_{\alpha} \vert$ 个渐近于
\begin{equation}\label{eq:ch05-tasym}
t \cong - \frac{\check c_{\alpha +1}}{\check c_\alpha} v^{(\check k_{\alpha+1} - \check k_{\alpha}
)}
\end{equation}
的根 $v(t)$。对这些根，$v(t)\to \infty$：其中 $t\to \infty$ 对应 $0\leq \alpha \leq \alpha_--1$，而 $t\to 0$ 对应 $\alpha_+\leq \alpha \leq n$。$t\to \infty$ 与 $t\to 0$ 处余下的根分别是 $\check g_0(v)$ 与 $\check g_{n+1}(v)$ 的根。这后一批根保持有限。此外，若 $\alpha_+>\alpha_-$，则在
\begin{equation}\label{eq:intermed}
\sum_{\alpha=\alpha_-}^{\alpha_+} \check c_{\alpha} t^{\alpha_+-\alpha} =0
\end{equation}
的根 $t\in \IC^\times$ 处，有 $\alpha_+-\alpha_-$ 个根 $v(t)$ 趋于无穷。一般地，这些将是单极点 $v(t) \sim \frac{\rho_\alpha}{t-t_\alpha}$。\footnote{在共形情形禁止这些根的 $U(1)_R$ 对称性，在这里是反常的。} 这样一来，牛顿多边形的边界就支配了 M5-膜的渐近形状。图 \ref{fig:brane-bending-2} 说明了这一点。

常数 $\check c_\alpha$ 的物理解释，现在可以从上面的图景推出。在同一 $v$ 值处测得的 M5-曲面相邻分支，其对应的 $t$ 值由下式相联系
\begin{equation}\label{eq:ch05-trunning}
\frac{t^{(\alpha+1)}}{t^{(\alpha)}} \cong \frac{ \check c_{\alpha+1} \check c_{\alpha
-1}}{\check c_\alpha^2} v^{\check b_\alpha}
\end{equation}
在弱耦合极限下我们已经看到，$\log\vert\frac{t^{(\alpha)}}{t^{(\alpha +1)}}\vert$ 应被解释为耦合常数的倒数。于是，对 $\check b_\alpha<0$ 我们看到它以尺度 $\vert v \vert$ 对数式跑动，而且紫外截断尺度是
\begin{equation}\label{eq:ch05-Lambda}
\Lambda^{-\check b_\alpha} =\frac{ \check c_{\alpha+1} \check c_{\alpha -1}}{\check c_\alpha^2}
\end{equation}
而对 $\check b_\alpha=0$，组合 $\frac{ \check c_{\alpha+1} \check c_{\alpha
-1}}{\check c_\alpha^2}$ 一如从前地编码紫外耦合。

现在我们准备描述相应的 Hitchin 系统。这是一个 $U(\check K)$ 系统，在 $t=0,\infty$ 处有奇点，在 \eqref{eq:intermed} 的根处也有奇点。在后一类奇点处 $\varphi_t$ 有正则奇点。在 $t\to
\infty$ 处，$\varphi_t$ 呈块对角形式。有一个正则奇点块，其留数由 $\check g_0(v)$ 的根给出。此外还有 $\alpha_-$ 个形如 \eqref{eq:ISPatInfty} 的不规则奇点块，其中应代入 $r= \check k_{\alpha +1}-\check k_{\alpha}$ 与 $\nu_i = - \check c_{\alpha}/\check c_{\alpha +1}$，$0\leq \alpha \leq \alpha_--1$。类似地，在 $t =\infty$ 处有一个由 $\check g_{n+1}(v)$ 的根支配的正则奇点块，以及 $n+1-\alpha_+$ 个形如 \eqref{eq:ISPatInfty} 的不规则奇点块，其中 $r= \check k_{\alpha}-\check k_{\alpha +1}$、$\nu_i = - \check c_{\alpha}/\check c_{\alpha +1}$，$\alpha_+\leq \alpha \leq n$。

导致不规则奇点的大质量极限，可以直接在 Hitchin 系统的层面上进行。例如，考虑 Higgs 场的 $\epsilon\to 0$ 极限：
\begin{equation}\label{eq:ch05-epslimit}
\frac{1}{t} \begin{pmatrix} 0 & 1/2 \\ - \rho/\epsilon & 0 \\
\end{pmatrix} +\frac{1}{t-\epsilon} \begin{pmatrix} 0 & 1/2 \\ \rho/\epsilon & 0 \\
\end{pmatrix}
\end{equation}
对固定的 $\epsilon$ 有两个正则奇点。然而当 $\epsilon\to 0$ 时，Higgs 场发展出我们所述类型的不规则奇点。

上述构造最简单的例子是 $n=1$、$\check K=2$。这是一个带 $d_1 = N_f$ 个基本味的 $U(2)$ 理论。若 $N_f=0$，则 $\check k_0=0, \check k_1 = 2, \check k_2 =0$，故 $\check b_1 =
-4$。曲线 \eqref{eq:Fhatcurvechck} 是 $\check c_0 t^2 + \check c_1 t
(v^2-u) + \check c_2=0$，在对 $t$ 作重新标度之后，该曲线成为标准的 $SU(2)$ SW 曲线
\begin{equation}\label{eq:Standard-SU2-curve}
v^2 - u - t - \frac{\Lambda^4}{t} = 0,
\end{equation}
从而 $t^2 \lambda^2 =\left( u + t + \frac{\Lambda^4}{t}\right)
(dt)^2$，其中 $\pm \lambda$ 是 $\varphi_t dt$ 的本征值。Higgs 场上的边界条件说：在规范等价意义下，
\begin{equation}\label{eq:ch05-bc0}
\varphi_t dt \to \frac{\Lambda^2}{\vert t \vert^{1/2}}
\begin{pmatrix} 0 & 1 \\ e^{-i \theta} & 0 \\ \end{pmatrix}
\frac{dt}{t} \qquad\qquad t\to 0
\end{equation}
\begin{equation}\label{eq:ch05-bcinf}
\varphi_t dt \to -  \vert t \vert^{1/2}
\begin{pmatrix} 0 & 1 \\ e^{i \theta} & 0 \\ \end{pmatrix}
\frac{dt}{t} \qquad\qquad t\to \infty
\end{equation}
其中 $t = \vert t \vert e^{i \theta}$。我们将在第 10 节% TODO 合稿改 \ref{section:SU2Examples}
非常详细地讨论这个例子。
